\chapter{The Five Phases of Socratic Prompting}

A formalized methodology for Socratic engagement with LLMs:

\section{Phase 1: The Raw Inquiry}

Begin with the messy, ambiguous question. Don't optimize for efficiency---optimize for \emph{thinking aloud}.

\begin{quote}
\emph{``I'm trying to understand why documentation gets ignored in tech teams. What might I be missing?''}
\end{quote}

\section{Phase 2: Reflective Clarification}

Interrogate the response. Don't accept it---challenge it.

\begin{quote}
\emph{``What assumptions underlie that answer?''} \\
\emph{``How would this change in a startup context vs. enterprise?''} \\
\emph{``Reframe this through a psychological lens rather than organizational.''}
\end{quote}

\section{Phase 3: Personal Synthesis}

Anchor abstract responses in your lived experience. This is where insight becomes actionable.

\begin{quote}
\emph{``That reminds me of when I rewrote a wiki page last month. Can you help me extract a pattern from that experience?''}
\end{quote}

\section{Phase 4: Operationalization}

Convert insight into structure. This is the translation from understanding to artifact.

\begin{quote}
\emph{``Based on this, write me a checklist I can use for future documentation work.''}
\end{quote}

\section{Phase 5: Recursive Loop}

Start again, from your now-evolved perspective.

\begin{quote}
\emph{``What question should I be asking next, now that I understand this?''}
\end{quote}

The cycle continues. Each iteration sharpens understanding. The end state is not an answer---it's clarity.

\begin{quotebox}
\textbf{Agentic Era Application:} These five phases now apply primarily during \emph{planning and review}, not execution. The phases are compressed but not eliminated. Master them to make approval decisions well.
\end{quotebox}
